\documentclass[10pt,landscape]{article}
\usepackage[margin=0.5in]{geometry}
\usepackage{multicol}
\usepackage{amsmath, amssymb, amsfonts}
\usepackage{physics}
\usepackage{enumitem}
\setlist{nosep}
\setlength{\columnsep}{0.25in}

\begin{document}
\begin{multicols}{2}

\section*{ENGR/PHYS 216 -- Complete Formula \& Concept Cheat Sheet}
\hrule
\vspace{0.3em}

%%%%%%%%%%%%%%%%%%%%%%%%%%%%%%%%%%%%%%%%%%%%%%%%%%%%%%%%%%%%
%%%%%%%%%%%%%%%%%%%% LECTURE 1 %%%%%%%%%%%%%%%%%%%%%%%%%%%%%
%%%%%%%%%%%%%%%%%%%%%%%%%%%%%%%%%%%%%%%%%%%%%%%%%%%%%%%%%%%%

\section*{Lecture 1: Vectors \& Motion in 1D/2D}

\subsection*{Vector Basics}

\textbf{Magnitude:}
\[
|\vec{A}| = \sqrt{A_x^2 + A_y^2 + A_z^2}
\]

\textbf{Unit Vector:}
\[
\hat{A} = \frac{\vec{A}}{|\vec{A}|}
\]

\textbf{Dot Product:}
\[
\vec{A}\cdot\vec{B} = A_xB_x + A_yB_y + A_zB_z = |\vec{A}||\vec{B}|\cos\theta
\]
Shortcut: zero $\Rightarrow$ perpendicular.

\textbf{Cross Product:}
\[
\vec{A}\times\vec{B} =
\begin{vmatrix}
\hat{i} & \hat{j} & \hat{k}\\
A_x & A_y & A_z\\
B_x & B_y & B_z
\end{vmatrix}
\]
Magnitude:
\[
|\vec{A}\times\vec{B}| = |\vec{A}||\vec{B}|\sin\theta
\]

Right-hand rule determines direction.

\subsection*{Kinematics (Constant Acceleration)}

\[
v = v_0 + at
\]
\[
x = x_0 + v_0 t + \frac{1}{2}at^2
\]
\[
v^2 = v_0^2 + 2a(x-x_0)
\]

Trick: If no time given, use third equation.

\subsection*{Projectile Motion}

Horizontal:
\[
x = v_0\cos\theta \, t
\]

Vertical:
\[
y = v_0\sin\theta \, t - \frac{1}{2}gt^2
\]

Time of flight:
\[
T = \frac{2v_0\sin\theta}{g}
\]

Range:
\[
R = \frac{v_0^2\sin(2\theta)}{g}
\]

Max height:
\[
H = \frac{v_0^2\sin^2\theta}{2g}
\]

Shortcut: Maximum range at $45^\circ$.

%%%%%%%%%%%%%%%%%%%%%%%%%%%%%%%%%%%%%%%%%%%%%%%%%%%%%%%%%%%%
%%%%%%%%%%%%%%%%%%%% LECTURE 2 %%%%%%%%%%%%%%%%%%%%%%%%%%%%%
%%%%%%%%%%%%%%%%%%%%%%%%%%%%%%%%%%%%%%%%%%%%%%%%%%%%%%%%%%%%

\section*{Lecture 2: Newton's Laws \& Forces}

\subsection*{Newton's Laws}

1. Inertia  
2. $\sum \vec{F} = m\vec{a}$  
3. Action = $-$Reaction  

Always draw FBD.

\subsection*{Common Forces}

Weight:
\[
W = mg
\]

Normal: perpendicular to surface.

Friction:
\[
f_s \le \mu_s N
\]
\[
f_k = \mu_k N
\]

Incline shortcuts:
\[
mg\sin\theta \text{ parallel}
\]
\[
mg\cos\theta \text{ perpendicular}
\]

Tension: same in massless rope.

%%%%%%%%%%%%%%%%%%%%%%%%%%%%%%%%%%%%%%%%%%%%%%%%%%%%%%%%%%%%
%%%%%%%%%%%%%%%%%%%% LECTURE 3 %%%%%%%%%%%%%%%%%%%%%%%%%%%%%
%%%%%%%%%%%%%%%%%%%%%%%%%%%%%%%%%%%%%%%%%%%%%%%%%%%%%%%%%%%%

\section*{Lecture 3: Work \& Energy}

\subsection*{Work}

\[
W = \vec{F}\cdot\vec{d} = Fd\cos\theta
\]

Area under $F$ vs $x$ graph.

\subsection*{Kinetic Energy}

\[
K = \frac{1}{2}mv^2
\]

\subsection*{Work-Energy Theorem}

\[
W_{net} = \Delta K
\]

\subsection*{Potential Energy}

Gravitational:
\[
U = mgh
\]

Spring:
\[
U = \frac{1}{2}kx^2
\]

\subsection*{Energy Conservation}

\[
K_i + U_i = K_f + U_f
\]

Shortcut: If conservative forces only, use energy.

%%%%%%%%%%%%%%%%%%%%%%%%%%%%%%%%%%%%%%%%%%%%%%%%%%%%%%%%%%%%
%%%%%%%%%%%%%%%%%%%% LECTURE 4 %%%%%%%%%%%%%%%%%%%%%%%%%%%%%
%%%%%%%%%%%%%%%%%%%%%%%%%%%%%%%%%%%%%%%%%%%%%%%%%%%%%%%%%%%%

\section*{Lecture 4: Momentum \& Collisions}

\subsection*{Momentum}

\[
\vec{p} = m\vec{v}
\]

\subsection*{Impulse}

\[
\vec{J} = \vec{F}\Delta t = \Delta \vec{p}
\]

Area under force-time graph.

\subsection*{Conservation of Momentum}

\[
\sum \vec{p}_i = \sum \vec{p}_f
\]

\subsection*{Elastic Collisions}

Momentum conserved  
Kinetic energy conserved  

1D shortcut:
\[
v_{1f} = \frac{(m_1-m_2)}{m_1+m_2}v_{1i} + \frac{2m_2}{m_1+m_2}v_{2i}
\]

%%%%%%%%%%%%%%%%%%%%%%%%%%%%%%%%%%%%%%%%%%%%%%%%%%%%%%%%%%%%
%%%%%%%%%%%%%%%%%%%% LECTURE 5 %%%%%%%%%%%%%%%%%%%%%%%%%%%%%
%%%%%%%%%%%%%%%%%%%%%%%%%%%%%%%%%%%%%%%%%%%%%%%%%%%%%%%%%%%%

\section*{Lecture 5: Rotational Motion}

\subsection*{Angular Kinematics}

\[
\omega = \omega_0 + \alpha t
\]
\[
\theta = \omega_0 t + \frac{1}{2}\alpha t^2
\]
\[
\omega^2 = \omega_0^2 + 2\alpha\theta
\]

Analogy: $x \leftrightarrow \theta$, $v \leftrightarrow \omega$, $a \leftrightarrow \alpha$

\subsection*{Tangential \& Radial Acceleration}

\[
a_t = r\alpha
\]
\[
a_r = \frac{v^2}{r} = r\omega^2
\]

\subsection*{Torque}

\[
\tau = rF\sin\theta
\]
\[
\tau = I\alpha
\]

\subsection*{Moment of Inertia (Common)}

Point mass:
\[
I = mr^2
\]

Solid disk:
\[
I = \frac{1}{2}MR^2
\]

Rod (center):
\[
I = \frac{1}{12}ML^2
\]

Parallel axis:
\[
I = I_{cm} + Md^2
\]

\subsection*{Rotational Energy}

\[
K = \frac{1}{2}I\omega^2
\]

%%%%%%%%%%%%%%%%%%%%%%%%%%%%%%%%%%%%%%%%%%%%%%%%%%%%%%%%%%%%
%%%%%%%%%%%%%%%%%%%% LECTURE 6 %%%%%%%%%%%%%%%%%%%%%%%%%%%%%
%%%%%%%%%%%%%%%%%%%%%%%%%%%%%%%%%%%%%%%%%%%%%%%%%%%%%%%%%%%%

\section*{Lecture 6: Angular Momentum}

\subsection*{Definition}

\[
\vec{L} = \vec{r}\times\vec{p}
\]

Rigid body:
\[
L = I\omega
\]

\subsection*{Conservation}

If $\sum\tau_{ext}=0$:
\[
L_i = L_f
\]

Trick: Skater pulling arms in increases $\omega$.

%%%%%%%%%%%%%%%%%%%%%%%%%%%%%%%%%%%%%%%%%%%%%%%%%%%%%%%%%%%%
%%%%%%%%%%%%%%%%%%%% LECTURE 7 %%%%%%%%%%%%%%%%%%%%%%%%%%%%%
%%%%%%%%%%%%%%%%%%%%%%%%%%%%%%%%%%%%%%%%%%%%%%%%%%%%%%%%%%%%

\section*{Lecture 7: Oscillations}

\subsection*{Hooke's Law}

\[
F = -kx
\]

\subsection*{Simple Harmonic Motion}

\[
x(t) = A\cos(\omega t + \phi)
\]

\[
\omega = \sqrt{\frac{k}{m}}
\]

\[
T = 2\pi\sqrt{\frac{m}{k}}
\]

Velocity:
\[
v = -A\omega\sin(\omega t + \phi)
\]

Acceleration:
\[
a = -\omega^2 x
\]

\subsection*{Pendulum (Small Angle)}

\[
\omega = \sqrt{\frac{g}{L}}
\]

\[
T = 2\pi\sqrt{\frac{L}{g}}
\]

%%%%%%%%%%%%%%%%%%%%%%%%%%%%%%%%%%%%%%%%%%%%%%%%%%%%%%%%%%%%

\section*{Universal Problem-Solving Strategy}

1. Draw diagram  
2. Choose coordinate system  
3. Write known equations  
4. Identify conserved quantities  
5. Solve algebraically before plugging numbers  

\vfill
\hrule
\centerline{All Core Definitions, Theorems, Formulas, Shortcuts Included}

\end{multicols}
\end{document}